% 前置内容: 前言
% 本文件在封面之后,目录之前由 AJbook.cls 的 \frontmatter 引入.

\chapter*{第一版前言(节选)}
\addcontentsline{toc}{chapter}{第一版前言(节选)}

数学是一门在今天以极其多样,缤纷多彩的面貌呈现自身的科学.因此, 当一个人想开始学习数学时, 自然而然地会产生这样的问题: 对于最初的学习步骤, 哪些领域应当居于优先地位.当然, 教育的方式也随时代而演变.在探索当前问题的过程中涌现出的基本见解, 已经持续不断地融入了教学之中.在这里, 数学所涉及的绝非复杂的细节, 而是一些经常重复出现,被证明十分重要, 并且本身就可以在初等水平上通过有意义的例子加以研究的基本模式.因此, 在大学数学教育中以分析学和线性代数的各一门入门课作为开端被证明是行之有效的, 这两门课通常是分开讲授的.两个领域相互补充, 并理想地包含着大量有趣的数学基本模式.是的, 人们完全可以说, 分析学与线性代数的方法对于几乎所有其他数学领域都是根本性的.

本书的正文内容相当于一门两个学期线性代数入门课程的量, 这门课我曾在明斯特大学多次讲授.大多数学生在决定攻读数学时, 对线性代数已经具备一定的预备知识, 例如关于向量运算或线性方程组的求解.相反, 根据全部经验, 他们不太熟悉线性代数那些使其具有如此普遍适用性的抽象概念构造.可以确定地说, 线性代数的这一抽象方面给许多学生带来了新的,不熟悉的困难.我决定不去掩饰这些困难, 而是有意识地把克服它们置于突出的位置.因此, 本书从一开头就十分注重清晰而系统,却又循序渐进地发展线性代数中习以为常的那些理论概念构造.凡是在后来的思考或推广中必须加以修正的临时性处理办法, 都尽可能加以避免.只有当某一主题领域的理论层面得到澄清之后, 才在充分发挥全部效用的前提下处理相应的计算方法.

不过, 像线性代数这样在很大程度上已远离其原始几何根源的理论, 如果不同时解释人们为什么以这种或那种方式行事,相应的策略是什么, 或者在某个定义背后所设想的主要应用情形有哪些, 那是很难消化的.为了涵盖这类问题, 在一门讲座中, 除了纯粹的教材内容之外, 还要在相当大的程度上阐明相应的动机背景.在教科书中, 这一成分往往只在很小的程度上得到体现, 因为否则就要求在逻辑严谨的数学步骤与或多或少带启发性和直观性的思考之间不断切换, 这自然不利于叙述的统一性与清晰性.而在本书中, 每一章都以一系列"预备说明"(Vorbemerkungen)开头, 其目的在于阐明相应章节的动机背景.从读者现有的知识水平出发, 将要处理的主要问题(包括相应的几何背景, 如果有的话)都得到说明; 此外, 可能的方法与求解策略,所得到解的形式, 以及与此相关的困难也得到讨论.建议在学习某一章的过程中根据需要多次查阅这些预备说明, 以便从中获得最大收益.不过我要明确指出, 这些引言在很大程度上属于合理性说明.因此它们并非以通常的数学严谨性撰写, 因而也不属于真正的教学内容.

本书内容的广度几乎没有特别之处.书中处理了向量空间及其线性映射,矩阵与线性方程组,行列式,多项式,特征值与标准形理论, 以及欧几里得与酉向量空间.其中有一节关于外积(以星号 $*$ 标出)的内容是可选的, 作为应用, 其中证明了行列式的一般拉普拉斯展开定理.按照本书的总体策略, 向量空间自同态标准形的推导是在主理想环上的模这一框架内进行的, 不过这样的模直到 6.3 节开头才被引入.如果读者希望只限于标准形理论的初等部分, 那么他可以在 6.1 节(特征值与特征向量)和 6.2 节(最小多项式与特征多项式)之后, 直接进入第 7 章关于欧几里得与酉向量空间的内容.

\clearpage

\chapter*{第五版前言}
\addcontentsline{toc}{chapter}{第五版前言}

近年来, 我收到了大量针对我的"线性代数"的意见与建议.例如, 汉堡的赖因哈德·迪斯特尔(Reinhard Diestel)在自己讲授线性代数课程的过程中对本书进行了批判性的审读, 并给我寄来了一份较长的修改建议清单.一些在过去多轮勘误中始终显得顽固的印刷错误, 由莫里茨·施密特(Moritz Schmitt)指出.衷心地感谢你们所有人的反馈!

我把这一切作为契机, 对本书为此次新版进行了彻底的修订, 不过并未改变其内容广度.因此, 本书一如既往地提供对线性代数的细致入门, 它面向一门两学期入门课程的要求.而新增的是一个新开发的习题训练(Aufgabentraining), 旨在引导独立处理习题.这里首先涉及对求解练习问题的某些基本原则与策略的一般性讨论.随后, 对每一节都示范性地详细给出若干精选问题的完整解答.我希望这个实质上超出了所谓"标准解答"展示方式的设计能够引起兴趣, 并对解题有所帮助!

明斯特, 2014 年 5 月

西格弗里德·博施

\clearpage
